\section*{M8: Necessary-and-sufficient characterization of $k=2$ saturation}
\label{sec:k2-iff}

\textbf{Motivation.} Two structurally different constructions
(\texttt{construct\_k2\_independent\_map\_saturation},
\texttt{construct\_k2\_repeated\_map\_saturation}, both in
\texttt{src/seion\_core/research\_v5/k2\_sharpness.py}) are already known to
attain $E_T^{\mathrm{proj}}=\rho M L_T$ exactly for every $\eta\in(0,1]$.
Prior work (\texttt{equality\_conditions.py}) audits each construction
against an informal checklist but never states or proves the underlying
\emph{iff} theorem those checklists are instances of, and its repeated-law
audit was stale (marked \texttt{OPEN\_IF\_REPEATED} even after the repeated-
law witness had already resolved it -- corrected alongside this theorem,
\S k2/equality\_conditions.py \texttt{K2-EQ-05}).

\textbf{Setup.} Binary $k=2$ chain: node \texttt{in} combines leaves
$a,b$; node \texttt{out} (the root) combines node \texttt{in}'s output and
leaf $d$. Laws $\mu_{\mathrm{in}},\mu_{\mathrm{out}}$ bilinear,
$\|\mu_{\mathrm{in}}\|_{\mathrm{op}}\le M$, $\|\mu_{\mathrm{out}}\|_{\mathrm
{op}}\le M$, orthogonal projector $P$ at node \texttt{in}, orthogonal
projector $P_{\mathrm{out}}$ at the root, closure-residual operator norm
$\rho_{\mathrm{in}}:=\|(I-P)\mu_{\mathrm{in}}\|_{\mathrm{op}}\le\rho$. No
restriction on ambient dimension, projector rank, or whether
$\mu_{\mathrm{in}}=\mu_{\mathrm{out}}$.

\textbf{Theorem.} Write $D:=(I-P)\mu_{\mathrm{in}}(a,b)$. Then
\[
  E_T^{\mathrm{proj}} \;=\; \big\| P_{\mathrm{out}}\,\mu_{\mathrm{out}}(D,d) \big\|,
\]
and $E_T^{\mathrm{proj}} = \rho M \|a\|\|b\|\|d\|$ (i.e.\ the universal
bound is attained with equality) \textbf{if and only if} all three of the
following hold simultaneously:
\begin{align}
  \text{(EQ1)}\quad & \|\mu_{\mathrm{out}}(D,d)\| = M\|D\|\|d\| \tag{outer law saturates $M$ on $(D,d)$}\\
  \text{(EQ2)}\quad & \|D\| = \rho\|a\|\|b\| \tag{closure map saturates $\rho$ on $(a,b)$}\\
  \text{(EQ3)}\quad & \mu_{\mathrm{out}}(D,d)\in\mathrm{Range}(P_{\mathrm{out}}) \tag{root projection loses nothing}
\end{align}

\textbf{Proof.} Bilinearity gives $F_{\mathrm{out}}-R_{\mathrm{out}}=
\mu_{\mathrm{out}}(F_{\mathrm{in}},d)-\mu_{\mathrm{out}}(R_{\mathrm{in}},d)=
\mu_{\mathrm{out}}(F_{\mathrm{in}}-R_{\mathrm{in}},d)=\mu_{\mathrm{out}}(D,d)$
(using $F_{\mathrm{in}}=\mu_{\mathrm{in}}(a,b)$, $R_{\mathrm{in}}=P\mu_{\mathrm
{in}}(a,b)$, so $F_{\mathrm{in}}-R_{\mathrm{in}}=D$), so
$E_T^{\mathrm{proj}}=\|P_{\mathrm{out}}\mu_{\mathrm{out}}(D,d)\|$ as claimed.
The universal bound is then obtained from the chain of non-negative
inequalities
\[
  E_T^{\mathrm{proj}}
  \;\overset{(\mathrm a)}{\le}\; \|\mu_{\mathrm{out}}(D,d)\|
  \;\overset{(\mathrm b)}{\le}\; M\|D\|\|d\|
  \;\overset{(\mathrm c)}{\le}\; M\rho\|a\|\|b\|\|d\|,
\]
where (a) is projection contractivity ($\|P_{\mathrm{out}}x\|\le\|x\|$),
(b) is the definition of $\|\mu_{\mathrm{out}}\|_{\mathrm{op}}\le M$, and
(c) is the definition of $\rho_{\mathrm{in}}\le\rho$. For real numbers
$0\le x\le y\le z\le w$, $x=w$ forces $x=y=z=w$ (a non-decreasing sequence
cannot return to its start value without every step being flat). Hence
$E_T^{\mathrm{proj}}=w$ iff (a), (b), (c) are individually equalities, which
are exactly EQ3, EQ1, EQ2 respectively (in that order). $\blacksquare$

\textbf{Both known witnesses are instances.} Checked mechanically
(\texttt{certify\_k2\_saturation}, \texttt{tests/research\_v5\_test\_k2\_
characterization.py}):
\begin{itemize}
  \item Independent-law witness: $D=\rho e_1$ (EQ2 tight since
  $(I-P)\mu_{\mathrm{in}}(x,y)=\rho x_0y_0e_1$ attains $\rho$ at $x=y=e_0$);
  $\mu_{\mathrm{out}}(D,d)=M\rho e_0$ (EQ1 tight, since $\mu_{\mathrm{out}}(x,y)
  =Mx_1y_0e_0$ attains $M$ at $x=e_1,y=e_0=d$); output already in
  $\mathrm{Range}(P_{\mathrm{out}})=\mathrm{span}(e_0)$ (EQ3 tight).
  \item Repeated-law witness: identical computation with
  $\mu_{\mathrm{in}}=\mu_{\mathrm{out}}=\mu$, $\mu(x,y)=Mx_1y_0e_0+\rho
  x_0y_0e_1$ -- same $D=\rho e_1$, same $\mu(D,d)=M\rho e_0$, all three
  conditions tight. Weight-sharing does not violate EQ1--EQ3 because a
  single bilinear map can have distinct norm-attaining input pairs for its
  structurally different terms ($x_0y_0$ term attains $\rho$ at $(e_0,e_0)$;
  $x_1y_0$ term attains $M$ at $(e_1,e_0)$) -- this is the precise mechanism
  behind the informal Section IX remark that ``weight sharing by itself
  does not destroy sharpness.''
\end{itemize}

\textbf{Sufficiency exercised predictively.} A third, independently
constructed instance not matching either prior witness's exact form
(ambient dimension 3, $\mathrm{rank}(P)=2$, $\mathrm{rank}(P_{\mathrm
{out}})=1$ on a different axis than $P$, $M=2$, $\rho=0.6$) is built purely
by satisfying EQ1--EQ3 directly and verified to saturate
(\texttt{test\_theorem\_predicts\_a\_novel\_saturating\_configuration\_
not\_in\_prior\_registries}), exercising the theorem's predictive content
rather than only re-deriving already-known cases.

\textbf{What this does and does not establish.} This is a complete
characterization of saturation \emph{for the binary $k=2$ chain with the
stated hypotheses} -- it reduces ``does this configuration saturate?'' to
three independently checkable local conditions, and explains (does not
merely tabulate) why both known witnesses work and why weight-sharing per
se is not an obstruction. It does \textbf{not} determine whether EQ1--EQ3
are simultaneously satisfiable within any further-restricted subclass (e.g.\
the gated-planar-rotation family, where they demonstrably are not fully
compatible -- that family's exact value is $\eta^2$, not $\rho M$, see M2
class B); that remains a separate, class-specific question. It also says
nothing about $k\ge3$: the argument is a two-term telescoping identity
specific to a single internal-node pair and does not extend mechanically
(see \texttt{k3/general\_upper\_envelope.tex} for the $k=3$ analysis, which
needed a materially different argument).

\textbf{Non-saturating example (necessity exercised negatively).} Take
$M=1$, $\rho=0.4$, dimension 3, $P=\mathrm{diag}(1,1,0)$ (rank 2),
$P_{\mathrm{out}}=\mathrm{diag}(1,0,0)$ (rank 1), $\mu_{\mathrm{in}}(x,y)=
0.3\,x\otimes y$-type random small-magnitude tensors (operator norm well
under $M$ by construction, not saturating), unit random leaves $a,b,d$.
Mechanically checked (\texttt{test\_non\_extremal\_configuration\_fails\_
at\_least\_one\_condition\_and\_does\_not\_saturate}): $E_T^{\mathrm{proj}}
\le\rho M\|a\|\|b\|\|d\|$ holds (never violated, consistent with the
universal bound), but the ratio is strictly below $1$ and at least one of
EQ1--EQ3 is strictly slack -- concretely, since $\|\mu_{\mathrm{in}}\|_{
\mathrm{op}}<\rho$ for this instance, EQ2 fails immediately (the closure
map cannot saturate $\rho$ if the whole law's magnitude is already below
$\rho$), which is enough on its own to block saturation regardless of the
other two conditions. This confirms the theorem's necessity direction is
not vacuous: generic, non-adversarial configurations genuinely fail to
reach the bound, and the failure is traceable to a specific named
condition rather than being a diffuse aggregate shortfall.

\textbf{Status.} \texttt{PROVED} (unconditional real-number chain-of-
inequalities argument, no floating-point-only evidence). Admissible class:
finite-dimensional real (identical argument holds verbatim over
$\mathbb{C}$ with the same norms) binary $k=2$ chain, arbitrary ambient
dimension and projector rank, arbitrary bilinear laws bounded by $M$ with
closure-residual $\le\rho$, independent or repeated. \texttt{approval\_
status: PENDING\_HUMAN\_REVIEW} per project policy -- this is a repository-
internal proof, not an externally reviewed one.
